\section{Rozwiązanie efektywne}
Aby zdefiniować rozwiązanie efektywne sprawdzamy sumaryczy obu funkcji celu dla wartości $\alpha$ i $beta$ od 1 do 10 \\ 
w tym celu napisany został kod który modyfikuje wartości $\alpha$ i $\beta$ w pętli 

\begin{lstlisting}[language=Python, caption={Wyznaczanie alpha i beta}]
    koszt_wyniki = []
    zadowolenie_wyniki = []
    wyniki_funkcji = []
    maksymalny_wynik = 0;
    for alpha in range(0, 11):
        beta = 10 - alpha + sys.float_info.epsilon
        alpha /= 10.0
        beta /= 10.0
        # Update objective function
        model.objective = alpha * koszt_dystrybucji
        - beta * suma_satysfakcji
    
        # Solve the model
        model.solve()
        print(alpha)
    
        # Record the wyniki
        calkowity_koszt = pulp.value(koszt_dystrybucji)
        calkowite_zadowolenie = pulp.value(suma_satysfakcji)
        wynik_funkcji = pulp.value(alpha * koszt_dystrybucji
        - beta * suma_satysfakcji)
        koszt_wyniki.append(calkowity_koszt)
        zadowolenie_wyniki.append(calkowite_zadowolenie)
        if wynik_funkcji > maksymalny_wynik:
            maksymalny_wynik = wynik_funkcji
        wyniki_funkcji.append(wynik_funkcji)
        
    print("maksymalny_wynik", maksymalny_wynik, wyniki_funkcji)

\end{lstlisting}
Najwyższy wynik zostął uzyskany dla $\alpha = 10$ i $\beta = 0$ i wynosił on \textbf{156.5}

\paragraph{Symulacja procesu podejmowania decyzji}
Przeprowadzone zostały symulacje dla 5 sytuacji:
\begin{enumerate}
    \item Tylko minimalizacja kosztów $\alpha = 1.0$ $\beta = 0.0$
    \item Priorytet na minimalizacji kosztów $\alpha = 0.8$ $\beta = 0.2$
    \item Równy podział $\alpha = 0.5$ $\beta = 0.5$
    \item Priorytet na satysfakcji klientów $\alpha = 0.2$ $\beta = 0.8$
    \item Tylko satysfakcja klientów $\alpha = 0.0$ $\beta = 1.0$
\end{enumerate}

Ponownie w celu przeprowadzenia symulacji napisano kod w pythonie 
\begin{lstlisting}[language=Python]
# Pseudo-code, assuming model setup as previously discussed
scemariusze = [
    (1.0, sys.float_info.epsilon),
    (0.8, 0.2),
    (0.5, 0.5),
    (0.2, 0.8),
    (sys.float_info.epsilon, 1.0)]
wyniki = []

for alpha, beta in scemariusze:
    model.objective = alpha * koszt_dystrybucji
        - beta * suma_satysfakcji

    model.solve()

    calkowity_koszt = pulp.value(koszt_dystrybucji)
    calkowite_zadowolenie = pulp.value(suma_satysfakcji)
    wyniki.append(
        (alpha, beta, calkowity_koszt, calkowite_zadowolenie)
        )

for idx, (alpha, beta, koszt, zadowolenie) in enumerate(wyniki):
    print(f"Krok {idx+1}:")
    print(f"  koszt: {alpha}, zadowolenie: {beta}")
    print(f" 
     Calkowity koszt: {koszt}, 
     zadowolenie klienta: {zadowolenie}\n")
\end{lstlisting}

